\section{The Path Relation}

An extrapolated value is still a path reading. It has not floated free of the
route by which it was obtained. The grammar must therefore ask two questions
that ordinary shorthand often collapses. What signed displacement has the path
carried? And what cost did the path spend to carry it? The first question is
about orientation and endpoints. The second is about the ledger of transitions
used to get there.

The difference appears sharply in the three-rung corridor. Begin at the mass
rung and return to the charge rung. One route is direct: mass to the returned
charge face. Another route is composed: mass to value, then value to returned
charge. If the residues compose correctly, both routes carry the same signed
displacement through the corridor. The grammar may identify that displacement.
It may say that the endpoint relation has returned the same oriented result.

In schematic form:
\[
  \Delta_{\pm}(\gamma_{12}\circ\gamma_{23})
    = \Delta_{\pm}(\gamma_{13}) .
\]
The symbol \(\Delta_{\pm}\) names the signed displacement. The composed path
\(\gamma_{12}\circ\gamma_{23}\) passes through an intermediate rung. The
direct path \(\gamma_{13}\) does not. The equality says only that the signed
endpoint displacement is the same. It does not say that the two paths have
spent the same cost.

Cost is a different reading:
\[
  C(\gamma_{12}\circ\gamma_{23})
    \quad\hbox{and}\quad
  C(\gamma_{13}) .
\]
The composed path spends two transitions. It carries the intermediate value
reading, restores the residue at that middle face, and then carries the next
transition. The direct path bypasses the middle face. Even when the signed
displacement identifies, the costs may differ. The grammar forbids the cost
ledger from being quotiented merely because the displacement has been
quotiented.

This is not a minor bookkeeping point. It is what keeps alpha finite. If the
grammar erased cost whenever endpoint displacement agreed, then every cheaper
and more expensive route with the same sign would become the same trial. The
bounded Richardson estimate would lose its bound, because the route by which
the bound was obtained would no longer be readable. A finite extrapolation
must know both what displacement was carried and what work was spent to carry
it.

The composed path carries an intermediate residue. That residue may settle
inside the value rung before the returned charge face is read. The direct
path does not expose that middle residue. It may arrive at the same signed
endpoint because the corridor relation composes. But the absence of an
intermediate exposure is itself a difference in cost. The grammar therefore
identifies the displacement while preserving the route.

This is the path version of the earlier charge/mass distinction. Charge
counts closure. Mass reads strain. A single closed loop can have both faces,
but the count is not the strain. In the path relation, signed displacement
counts the oriented change through the corridor. Cost reads the transitions,
settling, and residue restoration by which that change was obtained. They
belong to the same path family, but they are not the same reading.

The tail-latency example shows the point at a boundary. A message may have a
short displacement in presentation and still cost more because it must be
reconciled across many connected interfaces. The slowest admissible
reconciliation governs the update, not the shortest drawn route. The grammar
of path cost has the same form. A direct displacement may look short, but
the cost belongs to the reconciliation the route must actually spend.

The Sagnac example supplies the cyclic face. Two routes can traverse the
same boundary in opposite directions and return to comparable endpoints while
accumulating unequal refinement tallies. The difference is not local
distance. It is the holonomy of the bookkeeping rule. Here again the grammar
separates endpoint relation from path cost. A closed cycle can return to the
same reference while still charging a different cyclic debt.

For the alpha reading, this separation has a precise consequence. The
dimensionless result may be recovered from the signed charge relation after
the three-rung aperture and Richardson bound have done their work. But the
recovery is not complete until the path cost has been read as well. A direct
and a composed route that agree in signed displacement may support the same
\(\alpha\) estimate while certifying different costs of obtaining it.

The grammar therefore gives two verdicts, not one. On the displacement side,
it permits a quotient: the composed and direct readings identify when their
signed residue depth agrees. On the cost side, it carries the difference:
the route through the intermediate rung and the route that bypasses it are
not the same expenditure unless a separate cost bound forces them to be
indistinguishable. The quotient of displacement never erases the obligation
to read cost.

This also prevents a false contradiction. If two routes produce the same
signed displacement at different cost, the grammar has not failed. It has
learned more. It has read that the endpoint relation is invariant under the
composition, while the path expenditure is not. The invariant belongs to the
alpha estimate. The expenditure belongs to the finite process by which the
estimate was made.

Thus the path relation completes the chapter's core. Seed and residue make
the repeated trial one history. The three rungs make the aperture readable.
Depth doubling and sign-flipped error make the extrapolation bounded. The
path relation identifies signed displacement while preserving cost. Alpha
enters only after all four obligations have been carried.
